%
%- motivation for actual reading
%
We provide considerable detail on the characteristics and fundamental steps of method 
development of a novel geometrical unsplit VoF algorithm for arbitrary unstructured 
meshes. Our algorithm has shown to produce a bounded volume fraction field; moreover, 
it is mass conservative and robust, even though its advection part is based on 
a discrete Lagrangian flow map (velocities in mesh points), which has been 
reported to be unfeasible in literature so far. 

%
%- descriptive part
%
% detailed
The method's advection algorithm -- besides its use of an unique discrete 
Lagrangian flow map -- essentially comprises of a novel flux correction 
algorithm and narrow band computation: a fast iterative flux correction 
algorithm was implemented as well as an efficient narrow band computation 
based exclusively on the volume fraction field. This has been accomplished 
without introducing additional data structures and related algorithms 
required to update the narrow band as the interface moves. %
The key feature of the reconstruction algorithm is its accurate Node Averaged Gauss 
method used for calculating the phase fraction gradient, which results in lowest 
reconstruction errors reported so far for unstructured meshes.


The implementation of the geometrical intersection algorithms used both for 
reconstruction and advection is completely general, which makes the overall 
framework applicable to arbitrary cell types. %
However, the results show the algorithm to be more accurate on
unstructured hexahedral meshes than on tetrahedral meshes. In order to obtain
more accurate solutions on a tetrahedral mesh, higher order reconstruction
and/or advection algorithms will be required, which is to be expected. 
To increase the overall accuracy on hexahedral meshes, we have
extended the local dynamic Adaptive Mesh Refinement (AMR) engine in OpenFOAM
and implemented a geometrical mapping algorithm for the volume fraction field.
Results obtained using the adaptive geometrical VoF algorithm are promising and
are showing expected behaviour in absolute error values as well as error
convergence with a low computational overhead when compared to static mesh 
cases.  

The implementation of the geometrical library and the client simulation 
applications was done in a modular way using a policy based design in the C++ 
programming language, allowing for straightforward future extensions, if proven 
to be necessary. %
Our VoF algorithm is fully parallelized using the domain decomposition approach 
and is developed such that it is fully dimension agnostic following the design 
of the OpenFOAM library, which means that the transport automatically supports 
both two and three dimensional computation. The results are based on the majority 
of the standard established test cases for verification and validation of VoF 
algorithms and show that the new geometrical VoF algorithm delivers accurate and 
reliable results on unstructured meshes, even for large Courant numbers and 
complex velocity fields.

Our further research is directed towards the coupling of the geometrical VoF
algorithm with the pseudo-staggered flow solution algorithms in OpenFOAM on
arbitrary unstructured meshes. Another point of interest related to a
geometrically transported volume fraction field is the balancing of 
interfacial forces which will enable simulations of surface-tension
dominated flows in flow domains of arbitrary geometrical complexity. 
